% Cruzamento econ × espectro político
\begin{table}[htbp]
\centering
\SingleSpacing
\caption{Distribuição do espectro político por grupo (econ=0 vs.\ econ=1).}
\label{tab:crosstab_econ_esp}
\footnotesize
\begin{tabular}{lrrrrr}
\toprule
\textbf{Espectro} & \textbf{Públ.\ geral} & \textbf{\%} & \textbf{Economistas} & \textbf{\%} & \textbf{Total} \\
\midrule
Ext. esquerda & 6 & 4,5\% & 1 & 7,1\% & 7 \\
Esquerda & 37 & 28,0\% & 7 & 50,0\% & 44 \\
Centro & 21 & 15,9\% & 2 & 14,3\% & 23 \\
Direita & 59 & 44,7\% & 4 & 28,6\% & 63 \\
Ext. direita & 9 & 6,8\% & 0 & 0,0\% & 9 \\
\midrule
\textbf{Total} & \textbf{132} & 100\% & \textbf{14} & 100\% & \textbf{146} \\
\bottomrule
\multicolumn{6}{c}{\footnotesize Fonte: elaboração do autor.} \\
\multicolumn{6}{p{11cm}}{\footnotesize \textit{Nota:} Respondentes que se declararam ``Independente'' ou ``Sem opinião'' não possuem valor de espectro e foram excluídos desta tabela ($n_{\text{excl}} \approx$ correspondente à diferença entre $N=184$ e o total acima). Percentuais calculados dentro de cada coluna (grupo).} \\
\end{tabular}
\end{table}